\section{Atlas visual: todo lo que se midi\'o y se grafic\'o}
\label{sec:atlas}

Esta secci\'on existe por una raz\'on pr\'actica: el proyecto gener\'o
\textbf{715 archivos de imagen} entre gr\'aficos, esquem\'aticos, capturas y
diagramas. Ponerlos uno por p\'agina ser\'ian varios cientos de p\'aginas; no
ponerlos equivaldr\'ia a perder el registro de lo que efectivamente se midi\'o.

La soluci\'on adoptada son \textbf{hojas de contacto}: cada conjunto de figuras se
compone en una grilla rotulada, generada por el \emph{script}
\texttt{scripts/fig\_mosaico.py} directamente desde los directorios originales.
As\'i cada barrido, cada captura y cada an\'alisis queda \emph{visible} y con su
nombre de archivo, de modo que cualquiera de ellos puede recuperarse a tama\~no
completo cuando haga falta. Las figuras individualmente decisivas s\'i aparecen a
tama\~no grande, en las etapas correspondientes.

\begin{center}\footnotesize
\begin{tabularx}{\textwidth}{@{}l r X@{}}
\toprule
Conjunto & Figuras & Qu\'e contiene \\
\midrule
Tanda calibrada, tiempo & 65 & Barridos de osciloscopio de 4 canales, 21/07/2026 \\
Tanda calibrada, espectral & 65 & Preprocesamiento espectral de cada una de las anteriores \\
Campa\~na de calibraci\'on & 26 & Barridos durante el ajuste del potenci\'ometro \\
Verificaci\'on de calibraci\'on & 13 & Barridos de comprobaci\'on posterior \\
Tanda descartada 20/07 & 9 & La tanda con la resistencia BP incorrecta \\
Normalizados por etapa & 5 & Identificaci\'on BP, LP, compensador y cadena \\
C\'alculos r\'apidos (Python) & 10 & Exploraciones num\'ericas sueltas \\
Figuras del paper URUCON & 49 & Material del art\'iculo de auto-calibraci\'on \\
Assets de investigaci\'on & 51 & Figuras extra\'idas y diagramas de la revisi\'on bibliogr\'afica \\
Propuesta URUCOM & 36 & Im\'agenes de la propuesta \\
Historial (este documento) & 33 & Figuras generadas espec\'ificamente para esta ayuda-memoria \\
\midrule
\textbf{Total catalogado} & \textbf{362} & m\'as esquem\'aticos, PDF vectoriales y capturas sueltas \\
\bottomrule
\end{tabularx}
\end{center}

\subsection{La campa\~na de barridos de la cadena anal\'ogica}

Es el conjunto de datos experimentales m\'as grande del proyecto y el que sostiene
todas las afirmaciones sobre la respuesta real del \emph{front-end}. El m\'etodo
fue barrer la entrada con se\~nal de frecuencia creciente y registrar
\textbf{simult\'aneamente los cuatro puntos de la cadena} con el osciloscopio:

\begin{center}\footnotesize
\begin{tabular}{@{}cll@{}}
\toprule
Canal & Punto de la cadena & Color en las figuras \\
\midrule
CH1 & Salida del PGA (referencia de entrada) & Amarillo \\
CH2 & Salida del pasa-banda MFB & Azul \\
CH3 & Salida del compensador (resta) & Rojo \\
CH4 & Salida del pasa-bajos, entrada del ADC & Verde \\
\bottomrule
\end{tabular}
\end{center}

Cada barrido cubre una d\'ecada o menos, y el conjunto va desde
\SI{10}{\milli\hertz} hasta \SI{200}{\kilo\hertz}: siete \'ordenes de magnitud.
Las transferencias se calculan por cocientes entre canales, de modo que la
amplitud absoluta del generador no importa.

\begin{figure}[p]
\centering
\includegraphics[width=\textwidth]{mosaico_osciloscopio_tanda.png}
\caption{\textbf{Las 65 capturas de la tanda calibrada del 21/07/2026}, en el
dominio del tiempo. Cada celda son los cuatro canales de un barrido; el
r\'otulo indica la banda y el identificador de captura. Se aprecia a simple
vista el comportamiento cualitativo esperado: en las bandas bajas el canal rojo
(compensador) se hunde ---la muesca de cancelaci\'on--- mientras el verde (LP)
sigue al azul; en las bandas altas el verde cae y el azul se mantiene. Las
celdas con la se\~nal roja recortada corresponden a las capturas con saturaci\'on
de CH4 documentadas en el informe de la tanda.}
\label{fig:mosaico-osc-tanda}
\end{figure}

\begin{figure}[p]
\centering
\includegraphics[width=\textwidth]{mosaico_preproc_tanda.png}
\caption{\textbf{Preprocesamiento espectral de las mismas 65 capturas}. Para cada
barrido se estima la transferencia por ventanas, se cancelan los arm\'onicos de
red por m\'inimos cuadrados y se calcula la coherencia. Los puntos con coherencia
insuficiente se descartan autom\'aticamente; los intervalos lineales de las
capturas parcialmente saturadas se recuperan como submuestras en lugar de tirar
la captura entera (pol\'itica introducida en la versi\'on 8 de la cach\'e).}
\label{fig:mosaico-preproc-tanda}
\end{figure}

\begin{figure}[p]
\centering
\includegraphics[width=\textwidth]{mosaico_osc_cal0721.png}
\caption{\textbf{Campa\~na de calibraci\'on del 21/07/2026} (26 capturas): los
barridos tomados \emph{durante} el ajuste del potenci\'ometro de cancelaci\'on.
Registran el proceso de b\'usqueda, no el resultado final.}
\label{fig:mosaico-osc-cal}
\end{figure}

\begin{figure}[ht]
\centering
\includegraphics[width=\textwidth]{mosaico_osc_verif0721.png}
\caption{\textbf{Verificaci\'on posterior a la calibraci\'on} (13 capturas). Es la
tanda que se us\'o como principal para la identificaci\'on, por tener una \'unica
posici\'on declarada del potenci\'ometro.}
\label{fig:mosaico-osc-verif}
\end{figure}

\begin{figure}[ht]
\centering
\includegraphics[width=0.96\textwidth]{mosaico_osc_descartada0720.png}
\caption{\textbf{La tanda descartada del 20/07/2026} (9 capturas). Se archiv\'o al
descubrir que la rama BP ten\'ia montada una resistencia incorrecta
(\SI{43}{\kilo\ohm}), lo que hac\'ia que el numerador del compensador realizara un
amortiguamiento residual de $\zeta=83{,}661$ en lugar del objetivo
$\zeta_0=0{,}25$. Se conserva deliberadamente: es la evidencia de por qu\'e el
valor $83{,}661$ aparece en documentos intermedios y \textbf{no debe usarse} como
par\'ametro del dise\~no.}
\label{fig:mosaico-osc-descartada}
\end{figure}

\begin{figure}[ht]
\centering
\includegraphics[width=\textwidth]{mosaico_normalizados.png}
\caption{\textbf{Identificaci\'on por etapas, curvas normalizadas.} Cinco
transferencias: BP, compensador/PGA, LP, cadena PGA$\to$ADC y cadena
GEO$\to$ADC. Cada curva ---cruda, procesada, ajuste \texttt{tfest} y modelo
``PSoC High''--- lleva su m\'aximo a \SI{0}{\decibel} de forma independiente, para
comparar \emph{forma} sin que la ganancia absoluta domine el gr\'afico. Las
figuras sin normalizar siguen siendo la referencia para ganancia y para las
envolventes de tolerancia.}
\label{fig:mosaico-normalizados}
\end{figure}

\subsection{Resultados num\'ericos de la identificaci\'on}

El an\'alisis de la tanda calibrada produce un informe cuantitativo por etapa.
Conviene tenerlo aqu\'i porque es el resumen m\'as compacto del estado real del
\emph{front-end}.

\begin{table}[ht]
\centering\footnotesize
\caption{Identificaci\'on por etapas, tanda calibrada del 21/07/2026.
``Ajuste'' es el error RMS del modelo \texttt{tfest} contra los puntos medidos;
``vs.\ modelo'' es la discrepancia contra el modelo te\'orico con componentes
nominales. La distinci\'on es central: la primera columna dice si el modelo
\emph{describe} lo medido; la segunda, si el circuito \emph{se comporta} como se
dise\~n\'o.}
\label{tab:identificacion}
\begin{tabularx}{\textwidth}{@{}l r r r r X@{}}
\toprule
Etapa & Puntos & Coher.\ med. & Ajuste [\si{\decibel}] & vs.\ modelo [\si{\decibel}] & Estado \\
\midrule
BP (CH2/CH1)      & 85 & 0{,}968 & 6{,}49 & 7{,}57 & Parcial: polo inferior fuera de banda \\
COMP (CH3/CH1)    & 71 & 0{,}964 & \textbf{0{,}38} & \textbf{11{,}64} & Modelo bueno, circuito desviado \\
LP (CH4/CH3)      & 33 & 0{,}945 & 0{,}60 & \textbf{0{,}36} & Confiable en la banda \'util \\
PGA$\to$ADC (CH4/CH1) & 61 & \textbf{0{,}995} & \textbf{0{,}20} & 11{,}48 & Confiable; sub-s\'onica no observable \\
GEO$\to$ADC (calculada) & 61 & 0{,}995 & 0{,}20 & 11{,}48 & Compuesta con $H_{\text{geo}}$ nominal \\
\bottomrule
\end{tabularx}
\end{table}

\seleccionado{\textbf{El LP funciona.} En \SI{10.2046}{\hertz} las dos capturas
independientes dan \SI{13.192}{} y \SI{13.202}{\decibel}, frente a
\SI{13.978}{\decibel} del modelo. En toda la banda de ajuste
\SIrange{10.6}{984}{\hertz} el error RMS contra el modelo es de
\SI{0.435}{\decibel}. Es la \'unica etapa que cierra contra el dise\~no nominal
sin salvedades.}

\dato{\textbf{La muesca de cancelaci\'on es reproducible pero insuficiente.} El
modelo predice $CH3/CH1=\SI{-57.582}{\decibel}$ en \SI{10.2046}{\hertz}; la
medici\'on da aproximadamente \SI{-31.26}{\decibel}. Son
\textbf{\SI{26.33}{\decibel} menos de profundidad}. Dos capturas independientes
concuerdan entre s\'i dentro de \SI{0.42}{\decibel}
(\SI{-31.466}{} y \SI{-31.048}{\decibel}) y dos barridos distintos de
\SIrange{10}{100}{\hertz} dan \SI{-31.28}{} y \SI{-31.40}{\decibel} alrededor de
\SI{12}{\hertz}. Es decir: \emph{el circuito es repetible; lo que no alcanza es
la profundidad}.}

\dato{\textbf{Por qu\'e no alcanza: la resoluci\'on del ajuste.} Con componentes
nominales, el cursor objetivo para $\zeta_0=0{,}25$ est\'a en
\SI{634.539}{\ohm} y la cancelaci\'on exacta en \SI{632.556}{\ohm}: apenas
\textbf{\SI{1.98}{\ohm} de separaci\'on} sobre un recorrido de
\SI{2}{\kilo\ohm}, es decir una parte en mil del potenci\'ometro. Por la
profundidad medida de la muesca, la posici\'on efectiva estar\'ia en torno a
\SI{674}{\ohm}, unos \SI{39}{\ohm} por encima del objetivo. La estimaci\'on
compleja que combina BP y CH3 \textbf{no supera el criterio de confiabilidad en
ninguna captura}, porque errores sub-decibel en CH2 dominan el residuo cerca de
la cancelaci\'on; el programa lo informa expl\'icitamente en vez de presentar el
m\'inimo num\'erico como una resistencia f\'isica.}

\dato{\textbf{Banda extendida a \SI{10}{\milli\hertz}.} La captura
\texttt{XD\_ALL0017} (\SIrange{0.01}{0.05}{\hertz}) est\'a libre de saturaci\'on y
aporta 14 puntos entre \SI{10.95}{} y \SI{45.67}{\milli\hertz}. Coherencias
medianas: BP 0{,}995, compensador 0{,}988, LP 0{,}9998, cadena 0{,}988; errores
RMS contra el modelo de 0{,}96 / 2{,}11 / 0{,}33 / \SI{2.34}{\decibel}. El ajuste
global mantiene su l\'imite inferior en \SI{0.2}{\hertz} para que una \'unica
captura sub-s\'onica no domine la identificaci\'on, pero los puntos se muestran y
se comparan.}

\pendiente{\textbf{BP no cierra entre campa\~nas.} Localmente es muy coherente
(mediana 0{,}9996), pero las dos repeticiones nuevas difieren
\SI{0.75}{\decibel} en ganancia absoluta en \SI{10.2046}{\hertz} (0{,}732 y
$-0{,}022$ frente a \SI{0.772}{\decibel} nominal), y al fusionar todas las
campa\~nas hist\'oricas el ajuste degrada a \SI{6.485}{\decibel} RMS. Eso no es
baja coherencia dentro de cada barrido: es que \emph{las campa\~nas no coinciden
entre s\'i}, sobre todo en alta frecuencia. Y una diferencia peque\~na en BP
importa mucho, porque el compensador depende de una cancelaci\'on casi exacta.
LP, en cambio, s\'i permanece consistente al fusionar (\SI{0.596}{\decibel} RMS).}

\subsubsection{Higiene de datos aplicada}

Vale registrarlo porque forma parte del m\'etodo y explica los n\'umeros:

\begin{itemize}[itemsep=0.15em]
    \item La tanda principal tiene 30 capturas completas; de la serie \texttt{XD}
    se conservan 17 \'unicas (68 archivos CSV).
    \item \texttt{XD\_ALL0006} result\'o ser una \textbf{copia SHA-256 exacta} de
    \texttt{XD\_ALL0005} en los cuatro canales y se elimin\'o de la base. No
    quedan CSV duplicados por contenido.
    \item Desde la versi\'on 8 de la cach\'e, una saturaci\'on \textbf{ya no elimina
    la captura completa}: se excluye un ciclo local alrededor de cada recorte y
    los intervalos lineales continuos restantes se convierten en submuestras por
    etapa. En esta corrida se recuperaron parcialmente \textbf{22 combinaciones
    captura-etapa} de 9 capturas saturadas (5 BP, 5 compensador, 6 LP, 6 cadena).
    S\'olo 5 combinaciones quedaron descartadas, por conservar entre
    \SI{0.4}{} y \SI{2.3}{\percent} del barrido.
    \item Los \'ordenes de los modelos se ampliaron sumando dos polos por cada
    operacional comprendido en la transferencia, y el mismo incremento al
    numerador para conservar el grado relativo: BP 4/3, compensador 5/4, LP 4/2 y
    PGA$\to$ADC 11/8. GEO$\to$ADC se compone con $H_{\text{geo}}$ nominal y queda
    con 13 polos.
\end{itemize}

\seleccionado{\textbf{Conclusi\'on operativa de la campa\~na}, tal como qued\'o
asentada: las grabaciones son coherentes y suficientes para concluir que
\emph{el circuito es repetible}; el LP est\'a bien; el problema dominante sigue
siendo la calibraci\'on extremadamente fina de la cancelaci\'on. \textbf{No conviene
seguir con m\'as barridos anchos}: primero hay que recalibrar alrededor de
\SI{10.2046}{\hertz} con mayor sensibilidad vertical en CH3, y despu\'es repetir
un barrido estrecho.}

\subsection{Otros conjuntos de figuras del proyecto}

\begin{figure}[ht]
\centering
\includegraphics[width=\textwidth]{mosaico_urucom.png}
\caption{Material gr\'afico del art\'iculo URUCON sobre auto-calibraci\'on y de la
propuesta asociada (49 figuras): diagramas de la cadena anal\'ogica y digital,
flujo de calibraci\'on, salidas por etapa, respuesta del SM-24 y verificaciones de
la cadena.}
\label{fig:mosaico-urucom}
\end{figure}

\begin{figure}[ht]
\centering
\includegraphics[width=\textwidth]{mosaico_assets_investigacion.png}
\caption{Assets de la revisi\'on bibliogr\'afica (51 figuras): esquemas extra\'idos
y redibujados de los art\'iculos revisados, comparaciones de m\'etodos y diagramas
de apoyo del \emph{vault} de investigaci\'on.}
\label{fig:mosaico-assets}
\end{figure}

\begin{figure}[ht]
\centering
\includegraphics[width=\textwidth]{mosaico_propuesta.png}
\caption{Im\'agenes de la Propuesta URUCOM (36 figuras).}
\label{fig:mosaico-propuesta}
\end{figure}

\begin{figure}[ht]
\centering
\includegraphics[width=0.9\textwidth]{mosaico_calculos_rapidos.png}
\caption{C\'alculos r\'apidos en Python (10 figuras): exploraciones num\'ericas
sueltas hechas para decidir puntos concretos del dise\~no, conservadas porque
documentan el razonamiento aunque no formen parte de ning\'un \emph{pipeline}.}
\label{fig:mosaico-calculos}
\end{figure}
